\section{Conclusions}
\label{sec:conclusions}

The Rustikai project was developed with the goal of creating a distributed key-value store
capable of offering predictable and guaranteed performance in multi-tenant environments,
matching enterprise solutions like DynamoDB~\cite{Dynamo2022},
a requirement often overlooked in major existing open-source solutions such as Cassandra~\cite{Cassandra} and MongoDB~\cite{Mongo}.

While the developed system still leaves ample room for improvement,
it demonstrates how the proposed techniques can be effectively used to design
a distributed data store suitable for real-time performance guarantees.

The experimental results confirm the validity of the approach, as the system:
\begin{itemize}
    \item maintained controlled performance even under high load and in the presence of noisy clients,
    ensuring effective isolation between tenants and guaranteeing each a minimum share of throughput and an acceptable maximum latency;
    \item achieved performance gains due to scaling across multiple configurations, remaining within a small variance of the expected throughput.
\end{itemize}
Although testing was conducted on non-enterprise hardware due to resource constraints,
similar or improved performance is expected on high-end production systems.

\subsection{Limitations and Future Work}
\label{subsec:limitations}

Despite the promising results, the current implementation presents several limitations that open up
opportunities for future research and development.

The current system assumes a static set of peers,
whereas real-world deployments require the ability to dynamically join and leave the cluster.
Future work should include the development of a catch-up protocol for reintegrating nodes after outages or cluster membership changes,
as in Dynamo~\cite{Dynamo2007} and its derivatives.

Similarly, table throughput parameters are currently static and cannot be reconfigured at runtime.
A natural extension would involve dynamic reallocation of table throughput and key-space distribution,
enabling live adjustment of throughput guarantees.
Relaxing the uniform key-access assumption would further allow
an internal auto-scaling mechanism akin to Dynamo's partition-splitting strategy~\cite{Dynamo2007},
implemented using the dynamic table API\@.

While encryption is already supported by the underlying HTTP/2 server implementation,
most of the work required to enable TLS consists of exposing the configuration through the CLI\@.
Additional experimentation with HTTP/3~\cite{HTTP3} and QUIC~\cite{QUIC} could further improve latency and connection management.

Future work could also explore more sophisticated capacity management strategies across nodes,
particularly when tables have unbalanced throughput or storage requirements.
Enhancing the current throughput-limiting model to support separate limits for read and write operations
could be beneficial, drawing inspiration from DynamoDB's Read and Write Capacity Units (RCUs and WCUs)~\cite{Dynamo2022, DynamoWhitepaper}
or X-Stor's Request Units (RUs)~\cite{XStor}.

Several additional features could be built atop the existing system as middleware layers:
\begin{itemize}
    \item A fine-grained access control mechanism based on standard authentication and authorization frameworks such as OAuth or JWT~\cite{OAuth2, JWT};
    \item Built-in consistency options, including timestamp ordering or Dynamo-style vector clocks~\cite{Dynamo2007}, which could be provided transparently to the user;
    \item User-friendly table naming, implemented as a metadata layer stored within the system itself.
\end{itemize}

A parallel line of research, addressing similar concerns from a systems perspective, is that of CPU schedulers for real-time processes.
As developed by Lelli et al.~\cite{Deadline} for the Linux kernel with \texttt{SCHED\_DEADLINE},
the Earliest Deadline First (EDF) algorithm provides explicit latency guarantees for scheduled tasks.
Our implementation, through Rust's \texttt{tokio} runtime --- the de facto standard\footnote{
    \href{https://crates.io}{crates.io} reports over 416 million all-time downloads for \texttt{tokio},
    compared to approximately 70 million for the second most popular alternatives (\texttt{async-std} and its successor \texttt{smol}).
} for asynchronous programming in Rust --- uses a multi-threaded work-stealing scheduler
which was proven to be efficient in expected time, space, and communication~\cite{WorkStealingScheduler},
but provides no explicit latency control as some more advanced algorithms could.
On the same topic, both the operating system IO and packet scheduler could and/or should be enhanced by a similar approach,
but doing so would require a considerable effort.
